\scanpage{56}
14) \uline{Invariants de Galois et \unsure{thmes|thms} de commutation}

$X$ propre lisse sur $K$ de type fini

$H^{2i}(\bar X,\deleted{\mathbb Z}\mathbb Q_\ell(i))$ est de poids $0$, donc peut avoir des invariants non nuls sous Galois.

La conjecture de Tate s'énonce en disant que
\[
H^{2i}(\bar X,\deleted{\mathbb Z}\mathbb Q_\ell(i))^\pi
\xleftarrow{\sim}\unsure{\mathscr A|A|\mathcal A}^{i}(X/K)\otimes_{\mathbb Q}\mathbb Q_\ell
\]

On peut \unsure{le|la} préciser ainsi :

L'anneau des correspondances alg $/K$ \deleted{\unsure{en}}, \inserted{opérant en dim $i$} \deleted{\unsure{$C^i(X)$}} $\unsure{\mathscr U|\mathscr W|\mathcal U}^i(X/K)$ est un anneau semi-simple, et [via Hodge] et on veut que
\[
\begin{aligned}
&\mathscr U^i(X/K)\otimes_{\mathbb Q}\mathbb Q_\ell\ \text{soit le commutant de }\pi=\operatorname{Gal}(\bar K/K)\\
&\hspace{9em}\text{dans }\operatorname{End}_{\mathbb Q_\ell}(H^i(X,\mathbb Q_\ell(0))),
\end{aligned}
\]
\uline{et} l'alg. enveloppante de $\pi$ ds $\operatorname{End}_{\mathbb Q_\ell}(H^i)$ doit être le commutant de $\mathscr U^i(X/K)$ [donc \uline{l'action de $\pi$ est $1/2$ simple}]

Le deuxième énoncé $+$ injectivité de $\mathscr U^i(X/K)\otimes_{\mathbb Q}\mathbb Q_\ell \longrightarrow\operatorname{End}_{\mathbb Q_\ell}(H^i(X,\mathbb Q_\ell(0)))$ entraîne tout.
% archive sheet number: 55

